%% GENERATED by experiments/019-simb-multimodal/scripts/morphology_noise_ceiling.py
%% SOURCE: experiments/019-simb-multimodal/results/morphology_noise_ceiling.json, scalar_shortlist
%% Descriptions are verbatim from Ohya 2005 SI Table 2 "parameter statistics"
%%   $DATA_ROOT/torchcell-library/ohyaHighdimensionalLargescalePhenotyping2005/si/si6.pdf
%%   sha256 037a50927bc5b978ab38b207b264da87550feeeeaaf21bc311f65a5ef13e9638
\begin{table}[htbp]
\centering\footnotesize
%% Ragged-right paragraph columns: a narrow cell holding one long underscored token
%% cannot be justified, and justifying it emits an underfull hbox on every row.
\begin{tabular}{l >{\raggedright\arraybackslash}p{44mm} >{\raggedright\arraybackslash}p{72mm} r r r}
\toprule
feature & CalMorph label & SI description & ceiling & robust CV & score \\
\midrule
\file{A113} & \file{actin_n_ratio} & No actin patch found ratio & $0.873$ & $2.369$ & $2.068$ \\
\file{C124_C} & \file{Medium_bud_ratio} & Ratio of medium bud on stage C & $0.698$ & $0.886$ & $0.618$ \\
\file{A110_A1B} & \file{Actin_f_ratio} & Actin f ratio on stage A1B & $0.637$ & $0.939$ & $0.598$ \\
\file{C125_A1B} & \file{Large_bud_ratio} & Ratio of large bud on stage A1B & $0.784$ & $0.630$ & $0.494$ \\
\file{A105} & \file{actin_a_ratio} & Actin a ratio & $0.655$ & $0.617$ & $0.404$ \\
\file{D212} & \file{nuclear_B_ratio_to_nuclear_AA1BC_cells} & Ratio of B (Nuclear) to A, A1, B and C cells & $0.604$ & $0.642$ & $0.387$ \\
\file{D201} & \file{nuclear_B_ratio} & Ratio of B (Nuclear) & $0.565$ & $0.641$ & $0.362$ \\
\file{A109_A1B} & \file{Actin_e_ratio} & Actin e ratio on stage A1B & $0.476$ & $0.756$ & $0.360$ \\
\file{A103_A1B} & \file{Relative_distance_of_actin_patch_center_from_neck_in_mother} & Relative Distance of actin patch center from neck in mother cell on stage A1B & $0.812$ & $0.431$ & $0.350$ \\
\file{A103_C} & \file{Relative_distance_of_actin_patch_center_from_neck_in_mother} & Relative Distance of actin patch center from neck in mother cell on stage C & $0.534$ & $0.639$ & $0.342$ \\
\file{A110} & \file{actin_f_ratio} & Actin f ratio & $0.700$ & $0.468$ & $0.328$ \\
\file{D215} & \file{nuclear_B_ratio_to_nuclear_A1BC_cells} & Ratio of B (Nuclear) to A1, B and C cells & $0.480$ & $0.578$ & $0.277$ \\
\file{A8-2_C} & \file{Total_brightness_of_actin_region_in_bud} & Actin region brightness in bud on stage C & $0.594$ & $0.449$ & $0.267$ \\
\file{A8-2_A1B} & \file{Total_brightness_of_actin_region_in_bud} & Actin region brightness in bud on stage A1B & $0.595$ & $0.436$ & $0.259$ \\
\file{A121_A} & \file{Maximal_distance_between_patches} & Maximam actin patch length on stage A & $0.492$ & $0.517$ & $0.254$ \\
\file{A8-1_A} & \file{Actin_region_brightness} & Actin region brightness in mother cell on stage A & $0.562$ & $0.444$ & $0.250$ \\
\file{A120_A} & \file{Total_length_of_actin_patch_link} & Total length of actin patch link on stage A & $0.404$ & $0.613$ & $0.248$ \\
\file{D208} & \file{nuclear_B_ratio_to_budded_cells} & Ratio of B (Nuclear) to budded cells & $0.424$ & $0.575$ & $0.244$ \\
\file{A108} & \file{actin_d_iso_ratio} & Actin d (iso) ratio & $0.641$ & $0.375$ & $0.241$ \\
\file{C122} & \file{large_bud_ratio} & Ratio of large bud & $0.693$ & $0.342$ & $0.237$ \\
\bottomrule
\end{tabular}
\caption[]{The twenty best scalar morphology targets, ranked by ceiling multiplied by
robust CV. The list is dominated by bounded class ratios, actin localization class,
bud-size class and nuclear-stage class, which is what a deletion actually moves. Cell
size is the opposite trade and cell shape more so: cell size at the unbudded stage
\file{C11-1_A} has a ceiling of $0.928$ but a robust CV of $0.097$, ranking 101st, and
mother-cell roundness \file{C115_A} has a ceiling of $0.972$ with a robust CV of $0.024$,
ranking 245th of the 275 features with a defined rank key. That last id was described as a
whole-cell axis ratio in the previous revision; SI Table 2 calls it roundness of the mother
cell, and it is a shape parameter rather than a size one. Descriptions are quoted
verbatim from Ohya 2005 SI Table 2, source typography included, and the label column is
the same parameter's name in SI Table 1; the two tables word some parameters
differently, and both wordings are the paper's. In a CalMorph id the
leading letter names the stain the parameter is measured from, C the FITC-ConA cell-wall
image, A the rhodamine phalloidin actin image and D the DAPI nuclear image, and the
suffix names the nuclear stage the cells were binned into, \texttt{\_A}, \texttt{\_A1B}
or \texttt{\_C}, with no suffix meaning all stages pooled. The single-letter classes
inside a description, actin $a$ to $f$ and nuclear $A$ to $F$, are defined by the
drawings in SI Fig.~5B rather than by any text, so those rows rest on that figure.}
\label{tab:morph-scalar}
\end{table}
